% !TeX program = xelatex
\documentclass[12pt]{book}

% ----------------------------------------------------------
% Config

% A4
\usepackage[
  paperheight   = 297mm, 
  paperwidth    = 210mm,
  bindingoffset = 10mm,
  left          = 7.5mm,
  right         = 15mm,
  top           = 20mm,
  bottom        = 15mm,
  footskip      = 5mm,
]{geometry}

\usepackage[
  colorlinks,
  allcolors = black, 
]{hyperref}

\def\repoPath{/Users/phili/Code/Go/tsumego_workbooks/src}
\usepackage{\repoPath/config}

% ----------------------------------------------------------
% Text Config

\linespread{1.4}
% \setlength{\parindent}{0pt}
% \setlength{\parskip}{1.5em}

% ----------------------------------------------------------

\newcommand{\padfour}[1]{%
  \ifnum#1<10 000\fi%
  \ifnum#1<100 \ifnum#1>9 00\fi\fi%
  \ifnum#1<1000 \ifnum#1>99 0\fi\fi%
  #1%
}

% ----------------------------------------------------------

\begin{document}
  \workbookTitlePage[%
    title  = {%
      \vspace*{6cm}
      \fontsize{40}{40}\selectfont
      \textbf{수근대사전}
      \\
      \vspace{1cm}
      \scshape\fontsize{36}{36}\selectfont
      sugundaesajeon
    },
  ]
  
  \frontmatter

  \tableofcontents

  \chapter{Prefácio}
  
  \fontsize{17}{17}\selectfont

  Este livro de exercícios possui como fonte a Enciclopédia de Tesuji, da Nihon Kiin (日本棋院), a Associação de Go Japonesa. Na Coreia do Sul, intitulada como Sugeundaesajeon (수근대사전), diz-se que ela foi originalmente de autoria de Segoe Kensaku 9p e um de seus alunos mais famosos, um dos maiores ícones do jogo, Go Seigen 9p. Segoe Kensaku 9p possui seu próprio dicionário de tesujis, porém não está claro se é realmente o caso de ele ter tido alguma influência na enciclopédia atual.
  
  É bem provável que a influência seja concreta no entanto, pois Segoe Kensaku 9p fundou o dojo mais influente de sua época, treinando jogadores que viriam a revolucionar o Go, como Hashimoto Utaro 9p, Cho Hunhyun 9p e Go Seigen 9p.
  
  \section*{Como Usar Este Livro}
  
  Livros de exercícios não possuem respostas propositalmente. O inuito é desafiar o estudante a estabelecer a melhor variação para ambos os lados.
  
  % TODO: História da enciclopédia
  % TODO: Verdade que foi criado originalmente pelo Segoe Kensaku e Go Seigen?
  % Adicionar um diagrama demonstração de como se deve resolver um problema da coleçao.
  
  \mainmatter
  
  \setcounter{page}{1}

  \chapter{Problemas Locais}
  
  \pagebreak

  % Final problem for this section: 2476
  \foreach \i in {1,2,...,12}{
    \edef\paddedi{\padfour{\i}}%
    \problemDiagram[%
      problem size              = 7,
      path                      = \repoPath/sgf/sugeundaesajeon/problem-\paddedi.sgf,
      problem text              =,
      caption vertical margin   = 0.1cm,
      caption horizontal margin = -0.1cm,
      horizontal clip start     = 7,
      horizontal clip end       = 19,
      vertical clip start       = 1,
      vertical clip end         = 19,
    ]%
  }%
  \hspace*{-\fill}%
  
  \chapter{Problemas Globais}
  
  \pagebreak

  % Final problem for this section: 2634
  \foreach \i in {2477,2478,...,2482}{
    \edef\paddedi{\padfour{\i}}%
    \problemDiagram[%
      problem size              = 7.25,
      line width                = 0.475,
      path                      = \repoPath/sgf/sugeundaesajeon/problem-\paddedi.sgf,
      problem text              =,
      caption vertical margin   = 0.2cm,
      caption horizontal margin = -0.1cm,
      horizontal clip start     = 1,
      horizontal clip end       = 19,
      vertical clip start       = 1,
      vertical clip end         = 19,
    ]%
  }%
  \hspace*{-\fill}%
\end{document}